Interior design clients in Kyrgyzstan usually approve a project from 360\textdegree{} panoramas and renders. The designer chooses the camera, and wide-angle lenses and idealised lighting make rooms look larger and brighter than they will be. As a result, clients misjudge room size, light and how doors and furniture actually open, and these mistakes surface only during renovation.

This report presents \textit{JK3DStudio VR}, a tool that turns a finished 3ds Max interior project into an interactive virtual reality (VR) presentation. The client walks through the future apartment at 1:1 scale, opens doors by hand, switches light per room, changes the time of day and colour temperature, and is guided by an AI guide in Russian and English. The guide speaks only facts taken from a project passport that the designer approves.

The system has three parts: an Unreal Engine 5.8 editor plugin that prepares the scene semi-automatically from a naming standard, a VR application for PICO 4 Ultra, and an AI Gateway backend (FastAPI, PostgreSQL) for speech generation, fact checking and text-to-speech. Similar interactive tours abroad are produced manually and cost from USD 10{,}000; no comparable service for design projects was found in Kyrgyzstan. The novelty is fast preparation, honest scale and light, and a fact-grounded AI guide adapted to the local market. A working prototype of the VR application and a tested backend are presented at the first pre-defense.

\medskip\noindent\textbf{Keywords:} virtual reality, interior design, Unreal Engine, AI guide, large language models, Kyrgyzstan.
